\section{Titre III --- Temps de travail et congés}

\article{Durée du travail}\label{art:duree-travail}
I. --- La durée du travail est celle que fixe le contrat. Il n'existe pas de
durée légale du travail.

II. --- Le temps de travail se compte en minutes de travail effectif. La semaine
civile commence le lundi à zéro heure.
\lean{Contrat.dureeHebdo}
\lean{Planning}

\article[absolue]{Durée maximale}\label{art:duree-max}
La durée du travail, heures supplémentaires comprises, ne dépasse pas
quarante-huit heures par semaine civile en moyenne sur chaque période de
référence. Aucune dérogation individuelle n'est admise.
\lean{moyenne48h}
\theoreme{moyenne48h_independante}

\article[dispositive]{Période de référence}
La période de référence compte dix-sept semaines civiles. Par exception à
l'article~\ref{art:dispositives}, seul un accord d'entreprise peut la modifier,
dans la limite absolue de vingt-six semaines.
\lean{dispPeriodeReference}
\lean{periodeReference}
\theoreme{periodeReference_le_26}

\article[absolue]{Repos quotidien}
Deux journées de travail sont séparées par un repos d'au moins onze heures
consécutives. Une journée de travail s'étend sur treize heures au plus.
\lean{reposQuotidien}
\theoreme{travail_journee_le_13h}
\theoreme{reposQuotidien_independant}

\article[absolue]{Repos hebdomadaire}
Chaque semaine civile comporte un repos d'au moins trente-cinq heures
consécutives.
\lean{reposHebdomadaire}
\theoreme{reposHebdomadaire_independant}

\article[relative]{Pause}
Le salarié ne travaille jamais six heures d'affilée sans une pause d'au moins
vingt minutes.
\lean{pauses}
\lean{dispPause}
\theoreme{travail_continu_lt_6h}
\theoreme{pauses_independantes}

\article[absolue]{Travail de nuit}
I. --- Le travail de nuit est le travail accompli entre vingt et une heures
et six heures.

II. --- Est travailleur de nuit le salarié qui accomplit au moins trois heures
de travail de nuit, au moins deux nuits d'une même semaine civile.

III. --- Une journée de travail du travailleur de nuit compte huit heures de
travail effectif au plus.
\lean{travailleurDeNuit}
\lean{nuit8h}
\theoreme{travail_nuit_journee_le_8h}
\theoreme{nuit8h_independante}

\article[dispositive]{Heures supplémentaires}
I. --- Sont des heures supplémentaires les minutes de travail effectif qui
dépassent, dans une semaine civile, la durée hebdomadaire convenue au contrat.

II. --- Elles sont payées au taux horaire du contrat, majoré de dix pour cent.
Le taux horaire est le salaire mensuel divisé par la durée mensuelle convenue,
égale à la durée hebdomadaire multipliée par cinquante-deux et divisée par
douze.

III. --- La majoration ne peut être négative.
\lean{heuresSup}
\lean{paiementHS}
\lean{dispMajorationHS}
\theoreme{majorationHS_nonneg}
\theoreme{paiementHS_ge_normal}

\article[relative]{Congés payés}\label{art:conges}
Le salarié acquiert vingt jours ouvrés de congés payés par année de travail,
à raison de cinq tiers de jour par mois.
\lean{congesAcquis}
\lean{dispCongesAnnuels}
\theoreme{congesAcquis_annee}
\theoreme{conges_ge_legal}

\article[absolue]{Interdiction du rachat des congés}
Les congés payés ne peuvent être remplacés par une indemnité, sauf à la fin du
contrat pour les jours acquis et non pris.
\lean{indemniteCompensatrice}
\lean{dispRachatConges}
\theoreme{indemniteCompensatrice_fin}
\theoreme{pas_de_rachat_conges}
